% Fragmento integrable en VII: Relaciones estructurales entre las constantes físicas.
% Insertar después de desarrollos/normalizacion_conjunta.tex.
% Fuentes y cobertura íntegra: editorial/INSERCION_VII.md.
% Sin preámbulo ni macros; usa amsmath y los entornos ya definidos por VII.
\section{Gravedad, energía y autoinercia del sistema conjunto}
\label{hmtcuatro:gea:seccion}

La normalización conjunta permite formular una pregunta operatoria:
cómo intervienen la geometría, la materia y la memoria en una misma
realización. Se reúnen tres enlaces. El carácter positivo conserva
simultáneamente sus lecturas energética, másica y radial; la variación
del transporte conjunto produce corrientes geométricas e internas de
una misma energía; y la respuesta autoinercial se expresa mediante el
radio gravitatorio de ese carácter. Se demuestran la conservación de
estos enlaces bajo reducción de memoria y la propagación de las
restricciones de la acción acoplada durante su evolución regular.

El antecedente es el núcleo común del artículo. APP produce las dos
hojas con sus residuos, cocientes e incidencias; TRIT conserva régimen
y orientación: $+1$ corresponde a retorno y memoria, $0$ al umbral
presente y $-1$ a apertura. TPK selecciona, transporta, pliega, retorna
y conserva el estado enriquecido y su historia. La prolongación
$w_6\to w_{12}\to w_{18}\to w_{24}\to w_{30}\to R_{36}\to G_9$
mantiene prefijo, acarreo, supervivencia, ruta y frontera. La dinámica,
su holonomía nonádica y la publicación reducida tienen el orden
$U_t\to\Gamma_9\to K_{\rm ph}$: retorna la fase, avanza la memoria
y no se reinicia el estado. Las cinco construcciones consustanciales
de la estructura discreta del continuo no se separan al realizar sus lectores.

Las constantes y escalas empleadas son salidas de esta genealogía.
La cuadrirrelación $(\pi,\varphi,e,\alpha)$ conserva su origen común
y su orden constructivo interno; acción, velocidad constitutiva y
Barbero--Immirzi se reciben con sus lectores. El lenguaje operatorio
y variacional compone y reconoce esas salidas, sin introducir valores
metrológicos para seleccionar estados o coeficientes.

\subsection{Energía, masa y radio del mismo carácter}
\label{hmtcuatro:gea:caracter}

Se conserva la realización marcada de la sección anterior:
\[
\Omega=\{(j,k,q,u,v):j,k\in\mathbb Z/12\mathbb Z,\ 
k-j\in\{0,3,6,9\},\ 1\leq q\leq8,\ 
u<v,\ u,v\in\{1,2,4,5,7,8\}\}.
\]
Su cardinal $N_\Omega=12\cdot4\cdot8\binom62=5760$ cuenta incidencias,
no estados de materia ni lapsos. Para el inversor transportado
$B:V_{\rm in}\to V_{\rm out}$, con $B^*B=BB^*=I/4$, el proyector
$P=uu^*$ del modo común tiene $|u_i|^2=1/N_\Omega$. El archivo retiene
\[
C_0=(I-P)B,\qquad M_0=PB,\qquad C_0+M_0=B,\qquad\mathsf U=2B.
\]
Las imágenes de $C_0,M_0$ son ortogonales y $\mathsf U$ es unitario.
La expectativa marcada $\Delta_\Omega(A)=\sum_iP_iAP_i$,
$P_i=e_ie_i^*$, es única entre las expectativas bimodulares hacia
el álgebra diagonal que fijan los $P_i$: sobre una unidad matricial,
la bimodularidad exige
$\mathcal E(E_{ij})=P_i\mathcal E(E_{ij})P_j$, que es cero si $i\ne j$
y vale $E_{ii}$ si $i=j$.

De $C_0C_0^*=(I-P)/4$ y $P_iPP_i=P_i/N_\Omega$ resulta
\begin{equation}
\Delta_\Omega(C_0C_0^*)=r_\Omega I,\qquad
\Delta_\Omega(M_0M_0^*)=\frac{I}{4N_\Omega},\qquad
r_\Omega=\frac{5759}{23040}.
\label{hmtcuatro:gea:retorno}
\end{equation}
Ambas lecturas suman $I/4$. La regla constitutiva longitudinal aplica
linealmente el primer coeficiente a una sección de longitud:
\begin{equation}
L=\alpha^{16}r_\Omega L_*,\qquad D=L\mathsf U.
\label{hmtcuatro:gea:longitud}
\end{equation}
La norma local $\alpha^{16}$ y la sección $L_*$ son antecedentes
explícitos; tomar una raíz de intensidad sería otro lector.
La regla métrica radial es $\|R_Xv\|^2=\|XD^*v\|^2$.
Para $X=X^*>0$ en la fibra finita, la polarización y la raíz positiva
única dan $R_X^2=DX^2D^*=L^2\mathsf UX^2\mathsf U^*$ y
$R_X=LY$, con $Y=\mathsf UX\mathsf U^*>0$.
El espectro no se selecciona a partir de una respuesta gravitatoria deseada.

La misma sección $\hbar>0$ y el reloj conjugado $\omega L=c$
determinan $E_L=\hbar c/L$ y los lectores
\begin{equation}
H_X=E_LY,\qquad M_X=H_X/c^2,\qquad
R_X=LY,\qquad\Lambda_X=LY^{-1}.
\label{hmtcuatro:gea:lectores}
\end{equation}
La notación $H_X$ queda reservada a este lector del carácter.
El Hamiltoniano conjunto $\mathbb H$ introducido después tiene otro
dominio y no se identifica con $H_X$ por compartir el nombre de energía.

\begin{proposition}[Coeficiente radial común]
\label{hmtcuatro:gea:coeficiente}
En la identificación explícita de $R_X$ como radio gravitatorio reducido,
$R_X=GM_X/c^2$, existe un único coeficiente compatible:
\begin{equation}
G=\frac{c^3L^2}{\hbar},\qquad
R_X=\frac{G}{c^4}H_X=\frac{G}{c^2}M_X,\qquad
R_X\Lambda_X=L^2I,\quad H_X\Lambda_X=\hbar cI.
\label{hmtcuatro:gea:radial}
\end{equation}
\end{proposition}
\begin{proof}
Multiplicar los lectores prueba los dos productos. La identificación
radial equivale a $LY=G\hbar Y/(c^3L)$.
Cancelar el carácter invertible determina $G$; la sustitución inversa
verifica la igualdad para cada carácter positivo. Si dos coeficientes
la satisfacen, su diferencia multiplicada por $Y$ es cero y ambos coinciden.
\end{proof}

Las reglas longitudinal y métrica fijan la realización; la convención
de radio reducido fija el significado físico de $G$. El radio
Schwarzschild de la misma masa es $2R_X$. La igualdad conserva
superposiciones y elementos no diagonales en cualquier base transportada,
no sólo valores medios. Para una familia diferenciable y una conexión
de marco común $\mathcal A_t$, con $G,c$ fijos,
\begin{equation}
\mathcal D_tR_X=\frac{G}{c^4}\mathcal D_tH_X,\qquad
\mathcal D_tM_X=c^{-2}\mathcal D_tH_X,\qquad
\mathcal D_tO=\dot O+[\mathcal A_t,O].
\label{hmtcuatro:gea:derivada}
\end{equation}
Se obtiene diferenciando la identidad completa, incluido el marco.
No implica inmovilidad del estado.

Para modos positivos, o para operadores en factores tensoriales distintos,
\begin{equation}
m(y)=\frac{\hbar}{cL}y,\qquad
\frac{Gm_1m_2}{\hbar c}=y_1y_2,\qquad
\frac{R(y)}{\bar\lambda(y)}=y^2,\quad \bar\lambda(y)=L/y.
\label{hmtcuatro:gea:adimensional}
\end{equation}
La sustitución cancela las secciones dimensionales; los operadores en
factores distintos conmutan. El tamaño adimensional depende del carácter,
sin otro parámetro libre por especie. No se deduce de ello una ley
espacial de fuerza ni una evaluación numérica sin evaluar el carácter.

El calendario conserva $t_P=L/c$, $t_0=\pi L/(54c)$ y
$108t_0=2\pi t_P$. Por tanto $L=(54/\pi)\ell_0$, con $\ell_0=ct_0$;
las dos longitudes no se identifican. Entre secciones
$\hbar_b=\rho\hbar_a$, $\rho>0$, con $L,c,y_i$ fijos, se tiene
$G_b=G_a/\rho$, $m_{i,b}=\rho m_{i,a}$ y la cantidad
$Gm_1m_2/(\hbar c)$ permanece invariante. Se comparan secciones
conjuntas, no una variación temporal con las demás coordenadas
dimensionales mantenidas artificialmente fijas.

\subsection{Memoria estática, dinámica y dominio espectral}
\label{hmtcuatro:gea:memoria}

En dimensión finita, o con bloques acotados y coercivos, escribamos
\[
Y=\begin{pmatrix}A&F\\F^*&C\end{pmatrix}>0,\quad
Y_{\rm ef}=A-FC^{-1}F^*,\quad \eta=y+C^{-1}F^*b.
\]
Completar el cuadrado demuestra
\begin{equation}
\left\langle\binom b y,Y\binom b y\right\rangle
=\langle b,Y_{\rm ef}b\rangle+\langle\eta,C\eta\rangle.
\label{hmtcuatro:gea:schur-prueba}
\end{equation}
De aquí $Y_{\rm ef}>0$, la extensión minimizante
$y=-C^{-1}F^*b$ y la reconstrucción $y=\eta-C^{-1}F^*b$.
El residuo $\eta$ se conserva. Como
$\operatorname{Schur}(sY)=sY_{\rm ef}$ para $s>0$, se obtienen
$H_{X,\rm ef}=E_LY_{\rm ef}$ y $R_{X,\rm ef}=LY_{\rm ef}$.
Resolver $Y(x,y)=(f,0)$ muestra que el bloque de frontera de $Y^{-1}$
es $Y_{\rm ef}^{-1}$, luego $\Lambda_b=LY_{\rm ef}^{-1}$.
La compresión directa de $Y$ daría $A$: no es la misma operación.

Sea $\chi=L/E_L=G/c^4>0$ y $R_X=\chi H_X$.
En una descomposición compatible con los dominios, definimos
\[
\mathcal K_{H_X}(z)=H_{bb}-zI-V(H_{ii}-zI)^{-1}V^*.
\]
\begin{proposition}[Memoria dinámica completa]
\label{hmtcuatro:gea:resolvente-prop}
En el dominio resolvente del bloque interior,
\begin{equation}
\mathcal K_{R_X}(\chi z)=\chi\mathcal K_{H_X}(z).
\label{hmtcuatro:gea:resolvente}
\end{equation}
Si ambos complementos son invertibles, sus inversas llevan el factor
$\chi^{-1}$.
\end{proposition}
\begin{proof}
Cada bloque radial es $\chi$ veces el energético y
$(\chi H_{ii}-\chi zI)^{-1}=\chi^{-1}(H_{ii}-zI)^{-1}$.
La sustitución deja un factor $\chi$ en cada término, conservando el
orden. Invertir prueba la última afirmación.
\end{proof}
El parámetro energético $z$ se transporta a la longitud $\chi z$;
la identidad conserva cada orden de memoria espectral.
Para $|z|<\|H_{ii}^{-1}\|^{-1}$, la serie del resolvente da
\[
\mathcal K_{H_X}(z)=H_{\rm est}-zZ+O(z^2),\quad
H_{\rm est}=H_{bb}-VH_{ii}^{-1}V^*,\quad
Z=I+VH_{ii}^{-2}V^*\geq I.
\]
La positividad sigue de
$Z-I=(H_{ii}^{-1}V^*)^*(H_{ii}^{-1}V^*)$.
La extensión estática $(b,-H_{ii}^{-1}V^*b)$ tiene norma cuadrada
$\langle b,Zb\rangle$: energía y norma temporal proceden de la misma extensión.

Si $T=Z^{1/2}$ es acotado e invertible, el transporte es dual:
\begin{equation}
H_{X,c}=T^{-*}H_{X,\rm ef}T^{-1},\quad
R_{X,c}=T^{-*}R_{X,\rm ef}T^{-1},\quad
\Lambda_c=T\Lambda_bT^*.
\label{hmtcuatro:gea:dual}
\end{equation}
Multiplicar prueba $R_{X,c}=\chi H_{X,c}$,
$R_{X,c}\Lambda_c=L^2I$ y $H_{X,c}\Lambda_c=\hbar cI$,
sin conmutatividad de $T$ con el carácter.
Si $\psi=T(t)b$, se conserva $\dot\psi=\dot Tb+T\dot b$.
El término temporal de conexión del generador no se incorpora
retrospectivamente al carácter radial sin declarar su mapa.

\begin{proposition}[Extensión espectral y refinamiento compatible]
\label{hmtcuatro:gea:limite}
Sea $X$ autoadjunto, positivo e inyectivo, y $\mathsf U$ unitario.
Los lectores de \eqref{hmtcuatro:gea:lectores}, definidos por cálculo
espectral de $Y=\mathsf UX\mathsf U^*$, satisfacen
$R_X=\chi H_X$ en $\operatorname{Dom}(Y)$.
Sus productos con $\Lambda_X$ valen en $\operatorname{Dom}(Y^{-1})$
y tienen las prolongaciones acotadas de \eqref{hmtcuatro:gea:radial}.
Si inscripciones isométricas satisfacen $X'J=JX$,
$\mathsf U'J=K\mathsf U$ y conservan $L,E_L$, entonces
$R'_XK=KR_X$ y $H'_XK=KH_X$ en los dominios transportados.
\end{proposition}
\begin{proof}
Sobre $P_{\varepsilon,M}=\mathbf1_{[\varepsilon,M]}(Y)$, lectores e
inversos son acotados y las identidades son multiplicaciones de
funciones de $Y$. Para $\psi\in\operatorname{Dom}(Y)$,
\[
\|Y(P_{\varepsilon,M}-I)\psi\|^2
=\int_{(0,\infty)\setminus[\varepsilon,M]}
\lambda^2\,d\langle\psi,E_Y(\lambda)\psi\rangle\longrightarrow0.
\]
La inyectividad elimina masa espectral en cero y también
$P_{\varepsilon,M}\psi\to\psi$. Es convergencia en norma de grafo.
Si $\psi\in\operatorname{Dom}(Y^{-1})$, entonces
$Y^{-1}\psi\in\operatorname{Dom}(Y)$ y $YY^{-1}\psi=\psi$,
lo que prueba los productos y sus prolongaciones.
Sustituir los entrelazamientos en $R'_X=L\mathsf U'X'\mathsf U'^*$
prueba $R'_XK=KR_X$; la cuenta energética es igual.
\end{proof}
El sector nulo permanece separado. Se precisa así la extensión no
acotada de la normalización anterior. Una proyección ortogonal
arbitraria no queda habilitada para las fórmulas de bloques:
en operadores no acotados se requieren dominios compatibles o formas
cerradas con bloques controlados.

\subsection{Retroacción en el mismo transporte}
\label{hmtcuatro:gea:retroaccion}

El registro geométrico $x$, los enlaces internos $u$ y los campos $f$
proceden del mismo estado enriquecido. En una celularización orientada,
\begin{equation}
E(f,x,u)=\langle r,W(x,u)r\rangle,\quad r=D_Tf,\quad
(D_Tf)_a=f_{t(a)}-T_af_{s(a)},\quad
T_a=U_a^{\rm sp}(x)\otimes R_{\rm int}(u_a).
\label{hmtcuatro:gea:energia-memoria}
\end{equation}
$W=W^*>0$ conserva los bloques entre incidencias. Espín y
representación interna tienen factores distintos; la realización
quiral conserva sus bloques y el mapa Higgs--Yukawa equivariante.

\begin{proposition}[Diferencial de una misma energía]
\label{hmtcuatro:gea:variacion}
Con campos fijos y $q=Wr$,
\begin{align}
\delta E&=-2\operatorname{Re}\sum_a
\langle q_a,\delta T_af_{s(a)}\rangle+\langle r,\delta Wr\rangle,
\label{hmtcuatro:gea:diferencial}\\
\delta T_a&=\delta U_a^{\rm sp}\otimes R_{\rm int}(u_a)
+U_a^{\rm sp}\otimes\delta R_{\rm int}(u_a).\nonumber
\end{align}
\end{proposition}
\begin{proof}
Derivar la forma cuadrática da dos términos conjugados por
$W=W^*$ y el término $\langle r,\delta Wr\rangle$.
Usar $\delta r_a=-\delta T_af_{s(a)}$ prueba la primera fórmula;
la regla del producto prueba la segunda.
Ambas variaciones actúan sobre los mismos $r,q$, sin corrientes
elegidas independientemente.
\end{proof}

\subsubsection{Un espacio cuántico y un operador conjunto}
\label{hmtcuatro:gea:cuantica}

Se fija una celularización finita, un conjunto finito de registros
geométricos $\mathcal X$ y una Fock fermiónica finita $\mathcal F$.
Los enlaces recorren el producto compacto $\mathcal G^E$ recibido;
los dobletes $\Phi_v\in\mathbb C^2$ no tienen corte de amplitud.
Las medidas invariantes $d\mu_x=Z_x^{-1}e^{-S_x}d\mu_0$ tienen
densidades positivas acotadas con inversa acotada y no dependen de
$\Phi$ en esta carta. El espacio es
\[
\mathscr H=\bigoplus_{x\in\mathcal X}
L^2(\mathcal G^E\times\mathbb C^{2|V|},
\mu_x\otimes d\Phi;\mathcal F).
\]
Las identificaciones $A_xf=\sqrt{Z_x}e^{S_x/2}f$ son unitarias
desde la medida común y transportan el reloj a
$H_{\rm clk}^{\mu}=A(H_{\rm clk}\otimes I)A^*$.
El reloj mezcla registros geométricos del mismo espacio.
El operador se realiza mediante la forma de
\begin{equation}
\mathbb H=H_{\rm clk}^{\mu}+\hbar cK_{\rm gauge}
+d\Gamma(D_T^*WD_T)+H_{\Phi,\rm cin}+V_\Phi
+\widehat H_Y+\widehat H_{\rm tor}.
\label{hmtcuatro:gea:total}
\end{equation}
El mapa Yukawa es lineal en $\Phi,\bar\Phi$ y satisface
$Y(g\Phi)\rho_R(g)=\rho_L(g)Y(\Phi)$.
La contorsión eliminada no se conserva simultáneamente como variable
independiente en los términos base de esta carta efectiva.

La cinética de momentos escalares tiene forma
$\int\langle\nabla_\Phi\Psi,A(x,u)\nabla_\Phi\Psi\rangle$,
con $A$ acotada, uniformemente positiva e independiente de $\Phi$.
Se exige $A(x,u^g)=O_gA(x,u)O_g^{\mathsf T}$ para la acción real
$O_g$ de los dobletes; positividad sola no prueba esa covarianza.
La energía espacial escalar es un potencial cuadrático positivo
acotado por $C_{\rm grad}r^2$, $r^2=\sum_v|\Phi_v|^2$,
no un operador acotado sobre todo $L^2$.
Los volúmenes satisfacen
$0<w_{\min}\leq w_v(x)\leq w_{\max}$ y $\lambda_\Phi>0$.
Reloj, bloque gauge, transportes, $W$ y coeficientes torsionales
son acotados en la carta. El potencial conserva su valor absoluto:
\begin{equation}
V_\Phi=\sum_vw_v\lambda_\Phi(|\Phi_v|^2-v_0^2/2)^2
-\sum_vw_v\lambda_\Phi v_0^4/4.
\label{hmtcuatro:gea:potencial}
\end{equation}
La última contribución depende del volumen y no se sustrae.

\begin{theorem}[Forma cerrada y evolución conjunta]
\label{hmtcuatro:gea:forma}
La forma de \eqref{hmtcuatro:gea:total} es cerrada y acotada
inferiormente en
$\mathcal Q=\{\Psi\in\mathscr H:\nabla_\Phi\Psi\in L^2,\
r^2\Psi\in L^2\}$.
Determina un único operador autoadjunto asociado $\mathbb H$
y la evolución unitaria $e^{-it\mathbb H/\hbar}$.
\end{theorem}
\begin{proof}
Cauchy--Schwarz da
$\sum_vw_v\lambda_\Phi|\Phi_v|^4\geq ar^4$,
$a=\lambda_\Phi w_{\min}/|V|>0$.
Fock finita y Yukawa lineal dan
$\|\widehat H_Y(\Phi)\|\leq C_Yr$; la torsión está acotada.
Para $\varepsilon>0$,
\[
C_Yr\leq\varepsilon r^4+
\frac{3C_Y^{4/3}}{4^{4/3}\varepsilon^{1/3}},\qquad
br^2\leq\varepsilon r^4+\frac{b^2}{4\varepsilon}.
\]
La primera cota maximiza $C_Yr-\varepsilon r^4$; la segunda
completa el cuadrado en $r^2$. Aplicarlas con $\varepsilon=a/4$
a Yukawa y al cuadrático negativo da
\[
q(\Psi)\geq c_1\|\nabla_\Phi\Psi\|^2+
(a/2)\|r^2\Psi\|^2-C\|\Psi\|^2,\qquad c_1>0.
\]
La norma
$\|\Psi\|^2+\|\nabla_\Phi\Psi\|^2+\|r^2\Psi\|^2$
es completa por cierre de la derivada débil y la multiplicación.
Estas cotas y las superiores hacen equivalente la norma de forma,
tras sumar una constante, a esa norma completa.
El teorema de representación de formas cerradas produce el operador
asociado y su cálculo espectral da la evolución.
La unicidad se refiere a esta forma y dominio, no a cada dominio
formal alternativo.
\end{proof}

Covarianza de $D_T,W$, invariancia de medidas, identidad Yukawa y
contracción torsional interna conservan $q$ y $\mathcal Q$.
La unicidad del operador asociado implica que la representación
compacta conmuta con $\mathbb H$; su promedio de Haar reduce
el operador y su evolución. No es un promedio correctivo posterior.

Se retira el corte computacional con bloques completos de Peter--Weyl
y capas isotrópicas completas de Hermite.
Regularización y truncamiento suave aproximan derivadas y peso
$r^2$; las densidades acotadas preservan normas equivalentes.
Su unión es densa en $\mathcal Q$. Para las proyecciones ortogonales
$P_n$ en las medidas $\mu_x$ y los operadores de Galerkin
$\mathbb H_n$,
\[
(\mathbb H_n+\lambda)^{-1}P_nf\longrightarrow
(\mathbb H+\lambda)^{-1}f\qquad(\lambda>C).
\]
La ecuación coerciva de forma y su versión de Galerkin dan
ortogonalidad del error; continuidad y coercividad lo acotan por
el mejor error de aproximación, que tiende a cero.
Se retira el corte escalar y funcional de enlace de esta carta,
no el espacial, geométrico ni fermiónico.
Autoadjunción tampoco implica clase traza de $e^{-t\mathbb H}$.

\subsubsection{Intercambio y conservación}
\label{hmtcuatro:gea:intercambio}

Para $V=\sum_xP_x\otimes B_x$ y un elemento $h_{xy}$ del reloj,
multiplicar bloques da
\begin{equation}
[H_{\rm clk}\otimes I,V]_{xy}=h_{xy}(B_y-B_x).
\label{hmtcuatro:gea:conmutador}
\end{equation}
Si $h_{xy}\ne0$ y $B_y\ne B_x$, hay intercambio
geométrico--material. Si $\mathbb H=A+B+V$, para operadores
acotados o un núcleo común que justifique las derivadas,
\[
\frac d{dt}\langle A\rangle=\frac i\hbar\langle[\mathbb H,A]\rangle,
\quad
\frac d{dt}\langle B\rangle=\frac i\hbar\langle[\mathbb H,B]\rangle,
\quad
\frac d{dt}\langle V\rangle=\frac i\hbar\langle[\mathbb H,V]\rangle.
\]
La suma es cero por $[\mathbb H,\mathbb H]=0$.
Se conserva la energía incluyendo interacción, no cada sector
por separado. La geometría está en el espacio cuántico, no añadida
como fondo externo. Las historias ampliadas retienen memoria y retorno;
su refinamiento no equivale por definición al refinamiento espacial.

\subsection{La misma acción y sus corrientes}
\label{hmtcuatro:gea:accion}

La realización geométrica emplea cotetrada no degenerada $e$,
conexión de Lorentz $\omega$ y conexiones internas de la misma
materia. Sean $\Sigma^{IJ}=e^I\wedge e^J$,
$P_\gamma=\star+\gamma^{-1}I$,
$\gamma\in\mathbb R\setminus\{0\}$ y $\star^2=-I$ en bivectores.
Los acoplamientos permanecen fijos durante la variación.
Con $\kappa=8\pi G/c^4$, la misma acción es
\begin{equation}
S_{\rm tot}=\frac{\hbar}{16\pi L^2}
\int\Sigma_{IJ}\wedge(P_\gamma F_\omega)^{IJ}
-\frac{\Lambda\hbar}{8\pi L^2}\int\operatorname{vol}_e+S_m,
\label{hmtcuatro:gea:accion-total}
\end{equation}
donde $S_m=S_{\rm YM}+S_\Phi+S_{\rm Weyl}+S_Y$ contrae sus
índices con la misma $e$, que se varía.
Los coeficientes siguen de
$(2\kappa c)^{-1}=\hbar/(16\pi L^2)$.

Con $\delta_\omega S_m=-\tfrac12\int
\delta\omega^{IJ}\wedge\sigma_{IJ}$,
$\delta F=D_\omega\delta\omega$ y la integración por partes,
\begin{equation}
D_\omega(P_\gamma\Sigma)=\kappa c\sigma.
\label{hmtcuatro:gea:cartan}
\end{equation}
Se usan variaciones interiores o la completación de borde compatible.
Corriente de espín y fuente métrica son derivadas de la misma
materia, no dos entradas independientes.

\subsubsection{Eliminación torsional y términos cruzados}
\label{hmtcuatro:gea:torsion}

La inversa real es
$P_\gamma^{-1}=\gamma^2(-\star+\gamma^{-1}I)/(1+\gamma^2)$:
el producto de los dos polinomios en $\star$ antes de normalizar
vale $(1+\gamma^{-2})I$.
Denótese por $\mathcal Y^{IJ}=\kappa c(P_\gamma^{-1}\sigma)^{IJ}$
la forma bivectorial de grado tres, distinta del carácter $Y$.
Para $T^I=D_\omega e^I$,
$D_\omega\Sigma^{IJ}=T^I\wedge e^J-e^I\wedge T^J$.
Con el marco dual $E_I$, su inversa es
\[
\mathcal B^I=\sum_J\iota_{E_J}\mathcal Y^{IJ},\qquad
\vartheta=\tfrac14\sum_I\iota_{E_I}\mathcal B^I,\qquad
T^I=\mathcal B^I-e^I\wedge\vartheta.
\]
En efecto, si $t=\sum_J\iota_{E_J}T^J$, las contracciones dan
$\mathcal B^I=T^I+e^I\wedge t$ y después $4t$.
El mapa tiene núcleo nulo entre espacios de dimensión veinticuatro,
luego es un isomorfismo. Para
$\omega=\mathring\omega(e)+\mathfrak k$,
\[
T_{IJK}=\mathfrak k_{IKJ}-\mathfrak k_{IJK},\qquad
\mathfrak k_{IJK}=\tfrac12(T_{KIJ}-T_{IJK}-T_{JKI}).
\]
La combinación cíclica y la antisimetría en los dos primeros índices
prueban estas fórmulas.
Para materia afín en $\mathfrak k$, con corriente independiente
de ella, hay una única $\mathfrak k_*[e,\sigma]$ lineal en
la corriente total $\sigma=\sum_s\sigma_s$.

La expansión de curvatura conserva el borde
\[
S_\partial[e,\mathfrak k]=\frac1{2\kappa c}
\int_{\partial M}\Sigma_{IJ}\wedge(P_\gamma\mathfrak k)^{IJ}.
\]
La densidad interior dependiente de contorsión es
$\mathcal Q_{e,\gamma}(\mathfrak k)
-\tfrac12\mathfrak k^{IJ}\wedge\sigma_{IJ}$,
donde $\mathcal Q$ es homogénea de grado dos.
Su estacionariedad evaluada en
$\delta\mathfrak k=\mathfrak k_*$ da
$2\mathcal Q(\mathfrak k_*)=
\tfrac12\mathfrak k_*^{IJ}\wedge\sigma_{IJ}$.
Por tanto,
\begin{equation}
\mathcal L_{\rm spin}^{\rm ef}
=-\tfrac14\mathfrak k_*^{IJ}\wedge\sigma_{IJ}
=-\mathcal Q_{e,\gamma}(\mathfrak k_*).
\label{hmtcuatro:gea:spin}
\end{equation}
La linealidad conserva los cruces entre especies.
Se elimina la contorsión una vez; no se añaden simultáneamente
su interacción lineal eliminada y el cuártico efectivo.

En la especialización espinorial, el coeficiente recibido
$C_\gamma=3\kappa c\hbar^2\gamma^2/[16(1+\gamma^2)]$ da
\begin{equation}
\frac{C_\gamma}{\hbar}
=\frac{3\pi}{2}L^2\frac{\gamma^2}{1+\gamma^2}.
\label{hmtcuatro:gea:coeficiente-torsion}
\end{equation}
Basta sustituir $\kappa=8\pi G/c^4$ y $\hbar G=c^3L^2$,
sin cambiar la rama de $\gamma$.
En la Fock de la carta, la corriente total conserva la contracción
\[
\widehat Q_{{\rm tot},v}=\sum_{a=1}^4s_a
\{d\Gamma(\widetilde M_{a,v})^2-d\Gamma(\widetilde M_{a,v}^2)\},
\quad s=(-1,-1,-1,+1),
\]
\[
\widehat H_{\rm tor}
=-\sum_v\frac{N_vcC_\gamma}{\mathcal V_v}\widehat Q_{{\rm tot},v}.
\]
Las matrices $\widetilde M_{a,v}$ representan las componentes
de corriente. $N_v$ es el lapse, distinto de $N_\Omega$ y de
una ocupación. La identidad de orden normal retira la contracción
de una partícula, no los cruces entre especies.
La signatura permanece en $s_a$.

\subsubsection{Fuente métrica y Legendre de la misma acción}
\label{hmtcuatro:gea:legendre}

La eliminación estacionaria preserva la variación de las demás
variables: la regla de la cadena contiene un término en
$\delta\mathfrak k_*$ multiplicado por la ecuación de conexión,
que es cero. La acción reducida es
\[
S_{\rm red}=\frac1{2\kappa c}\int(R[g]-2\Lambda)\operatorname{vol}_e
+S_m^{\rm LC}+S_{\rm spin}^{\rm ef}+S_\partial.
\]
Definimos la fuente de acción por
$\delta_g(S_m^{\rm LC}+S_{\rm spin}^{\rm ef})
=-\tfrac12\int\mathcal T^S_{\mu\nu}\delta g^{\mu\nu}
\operatorname{vol}_e$.
La variación gravitatoria interior produce
\begin{equation}
G_{\mu\nu}+\Lambda g_{\mu\nu}
=\kappa c\mathcal T^S_{\mu\nu}
=\kappa(T^{\rm LC}_{\mu\nu}+U^{\rm spin}_{\mu\nu}),
\quad T^{\rm LC}+U^{\rm spin}=c\mathcal T^S.
\label{hmtcuatro:gea:metrica}
\end{equation}
Se incluye la variación de volumen del término constante del
potencial. El factor $c$ corresponde a la cotetrada en la recta
de longitud; no se mezclan las fuentes de acción y físicas.

Con $a=(2\kappa c)^{-1}$ y
$\mathcal B_{IJ}=(P_\gamma\Sigma)_{IJ}$, el potencial geométrico
es $\theta_{\rm geom}=a\mathcal B_{IJ}\wedge\delta\omega^{IJ}$.
Descomponer $\omega=\omega_0dt+\underline\omega$ y
$F_{0a}=\dot\omega_a-D_a\omega_0$ identifica el término cinético;
integrar por partes conserva Gauss y borde.
Con $e_0=Nn+N^ae_a$ y la Legendre de la misma materia,
\begin{equation}
S_{\rm tot}=\int dt\,[\Theta_{\rm red}(\dot z)
-\mathcal C[N]-\mathcal V[N^a]-\mathcal G_L[\omega_0]
-\mathcal G_{\rm int}[A_0]]+S_\partial.
\label{hmtcuatro:gea:canonica}
\end{equation}
Lapse, shift y calibre mantienen sus direcciones degeneradas.
Las restricciones de segunda clase se reducen antes de usar
$\Theta_{\rm red}$ o se conservan expresamente.
Los espinores de primer orden no reciben una inversa ficticia
de velocidades.

Para contorsión y su momento, las restricciones son
$\pi_{\mathfrak k}=0$ y
$\partial\mathcal Q/\partial\mathfrak k-\sigma/2=0$.
Su matriz tiene bloques
$\left(\begin{smallmatrix}0&-M\\M^{\mathsf T}&B\end{smallmatrix}\right)$,
con $M$ invertible por la reconstrucción anterior; por tanto es
invertible. El corchete de Dirac efectúa esa eliminación y conserva
separadamente las restricciones fermiónicas graduadas.
En una dirección bosónica regular,
\[
H_\Phi=N\left(\frac{p^2}{2w}+wV+H_{\rm grad}\right),\quad
p=\frac wN D_t\phi,\quad
L_\Phi=\frac{w}{2N}|D_t\phi|^2-NwV-NH_{\rm grad}.
\]
El volumen $w=\sqrt{\det q}$ no se congela.
Para $K_{ij}=(\dot q_{ij}-\mathcal L_{N^a}q_{ij})/(2N)$,
\[
\Pi^{ij}=\frac{\sqrt q}{2\kappa c}(K^{ij}-q^{ij}K),\qquad
K_{ij}=\frac{2\kappa c}{\sqrt q}
(\Pi_{ij}-\tfrac12q_{ij}\Pi).
\]
La traza y la parte sin traza prueban la inversa.
Si el momento total contiene una contribución material de marco,
se usa $\Pi=p-p_m$, conservando $p_m$.
Así se recupera la misma acción y su fuente. La Legendre no
demuestra por sí sola que una cuantización arbitraria de
\eqref{hmtcuatro:gea:canonica} coincida con
\eqref{hmtcuatro:gea:total}.

\subsection{Propagación de restricciones con la fuente total}
\label{hmtcuatro:gea:restricciones}

Sea $E_{\mu\nu}=G_{\mu\nu}+\Lambda g_{\mu\nu}
-\kappa c\mathcal T^S_{\mu\nu}$.
Sobre las ecuaciones materiales, la invariancia por difeomorfismos
del funcional reducido implica
$\nabla^\mu\mathcal T^S_{\mu\nu}=0$:
una variación interior generada por $\xi$ deja la integral de
$\mathcal T^S_{\mu\nu}\nabla^\mu\xi^\nu$, mientras las ecuaciones
materiales anulan los demás términos. Integrar por partes y variar
$\xi$ prueba la divergencia nula, antes de imponer $E=0$.
Bianchi y los acoplamientos fijos dan $\nabla_\mu E^{\mu\nu}=0$.

\begin{proposition}[Conservación de los datos de restricción]
\label{hmtcuatro:gea:propagacion}
Una solución regular de las ecuaciones materiales y evolutivas
métricas conserva las restricciones satisfechas inicialmente
durante su intervalo de existencia regular.
\end{proposition}
\begin{proof}
En coordenadas gaussianas
$g=-d\tau^2+q_{ij}(\tau,x)dx^idx^j$ se imponen $E_{ij}=0$ y
se escribe
$C=E_{00}$, $C_i=E_{0i}$,
$K_{ij}=\dot q_{ij}/2$ y $\mathcal K=q^{ij}K_{ij}$.
Se tiene $E^{00}=C$, $E^{0i}=-C^i$, $E^{ij}=0$,
$\Gamma^i{}_{0j}=K^i{}_j$ y
$\Gamma^\mu{}_{\mu0}=\mathcal K$.
La componente temporal de la divergencia y la espacial dan
\[
\partial_\tau C-D_iC^i+\mathcal KC=0,\qquad
\partial_\tau C^i+\mathcal KC^i+2K^i{}_jC^j=0.
\]
Al bajar el índice, $\dot q_{ij}=2K_{ij}$ cancela el último
término espacial. Quedan exactamente
\begin{equation}
\partial_\tau C-D_iC^i+\mathcal KC=0,\qquad
\partial_\tau C_i+\mathcal KC_i=0.
\label{hmtcuatro:gea:propagacion-sistema}
\end{equation}
La segunda ecuación homogénea mantiene $C_i=0$ y la primera
queda $\dot C+\mathcal KC=0$, conservando $C=0$.
$q$ y $\mathcal K$ son los de la solución acoplada, no un fondo prescrito.
\end{proof}
La carta gaussiana prueba localmente una afirmación tensorial,
sin congelar $q^{-1}$ ni sustituir la fuente.
No se afirma existencia global sin singularidades para cada dato.
La propagación variacional tampoco es por sí sola una identidad
de conmutadores cuánticos para deformaciones normales arbitrarias.

\subsection{Frontera y operador autoinercial}
\label{hmtcuatro:gea:autoinercia}

Sea $b:\mathcal X_\infty^{\rm enr}\to B_{\rm ef}$,
$B_{\rm ef}=\operatorname{im}(b)$.
Una respuesta $I_v$ factoriza unívocamente como
$I_v=\overline I_v\circ b$ si y sólo si es constante en las
fibras de $b$. La necesidad es inmediata; para la suficiencia,
$\overline I_v(\beta)=I_v(x)$ con $b(x)=\beta$ es independiente
del representante. La imagen efectiva da unicidad.
La dependencia machiana fuerte añade estados con igual restricción
local y respuestas distintas.

La bisagra del mismo estado produce
$r_{\rm av}=\pi/\varphi^2$ y $r_{\rm av}^J=\varphi/\pi^2$.
En el sector de amplitud común,
$\mathcal L_{\rm ar}=\pi\mu_{\rm ar}\mathcal M(x)$ y
$\mathcal L_{\rm vol}=\mu_{\rm vol}\mathcal M(x)$,
$\mathcal M(x)>0$, con
$\mu_{\rm ar}/\mu_{\rm vol}=\varphi^{-2}$.
Dividir por $\mathcal L_{\rm vol}$ da lectores del mismo tipo
$A_\partial=r_{\rm av}$ y $V_\partial=1$.
El cociente cancela la amplitud, no la historia conservada en el archivo.

Para una orientación $a$,
$q_a(x)=h_a(x)/(108t_0(x))=\hbar_a(x)/t_P(x)$
desciende a $Q:B_{\rm ef}\to\mathcal L_E^+$
si y sólo si $b(x)=b(x')$ implica
$h_a(x)/t_0(x)=h_a(x')/t_0(x')$.
Es el criterio de factorización aplicado al reloj.
Lo satisfacen las cartas fijas y las variables cuyo cociente
desciende. En la carta radial $Q=\hbar_ac/L$.
Con energías de esa misma recta, $E_1E_2/Q^2$ es adimensional
e independiente de base.

Para proyecciones ortogonales complementarias,
\begin{align}
w_A&=\frac{A_\partial r_{\rm av}}
{A_\partial r_{\rm av}+V_\partial r_{\rm av}^J}
=\frac{\pi^4}{\pi^4+\varphi^5},\qquad w_V=1-w_A,
\label{hmtcuatro:gea:pesos}\\
W_{\rm av}&=w_AP_A+w_VP_V,\qquad
\overline I_v=\frac{E_1E_2}{Q^2}W_{\rm av}.
\label{hmtcuatro:gea:respuesta}
\end{align}
Multiplicar numerador y denominador de
$r_{\rm av}^2/(r_{\rm av}^2+r_{\rm av}^J)$ por
$\varphi^4\pi^2$ prueba la evaluación.
Se conserva la diferencia entre lector areal y ponderación:
la marca regional se genera una vez.
$0<w_A,w_V<1$ hace positivo e invertible a $W_{\rm av}$ y
\[
\langle z,\overline I_vz\rangle
=\frac{E_1E_2}{Q^2}
(w_A\|P_Az\|^2+w_V\|P_Vz\|^2)>0\quad(z\ne0).
\]
Intercambiar completamente
$(A_\partial,r_{\rm av},P_A)$ y
$(V_\partial,r_{\rm av}^J,P_V)$ permuta los sumandos.
Un transporte unitario conjuga la respuesta y conserva su espectro;
implementar el intercambio dentro de un espacio exige iguales
dimensiones sectoriales.
Una reescala común de $A_\partial,V_\partial$ no cambia los pesos
y no prueba dependencia global.
Con energías y pesos fijos, $Q\mapsto qQ$ cambia la respuesta
por $q^{-2}$; el testigo fuerte conserva además la misma
restricción local.

\begin{theorem}[Autoinercia y radio del mismo carácter]
\label{hmtcuatro:gea:autoinercial-radial}
Con el mismo reloj y orientación,
\begin{equation}
\overline I_v=\frac{G_aE_1E_2}{\hbar_ac^5}W_{\rm av}
=\frac{G_am_1m_2}{\hbar_ac}W_{\rm av}.
\label{hmtcuatro:gea:mach-gravedad}
\end{equation}
Para una sonda $y_p>0$, $E_p=Qy_p$ y
$\bar\lambda_p=L/y_p$, la extensión espectral satisface
\begin{equation}
\widehat I_{p,Y}=y_pY\otimes W_{\rm av}
=\frac{E_p}{Q^2}H_X\otimes W_{\rm av}
=\frac{R_X}{\bar\lambda_p}\otimes W_{\rm av}
=\frac{G_am_p}{\hbar_ac}M_X\otimes W_{\rm av}.
\label{hmtcuatro:gea:identidad-autoinercial}
\end{equation}
\end{theorem}
\begin{proof}
$Q^2=\hbar_a^2c^2/L^2$ y $\hbar_aG_a=c^3L^2$ dan
$Q^{-2}=G_a/(\hbar_ac^5)$; después se usa $E_i=m_ic^2$.
Para $Y=\sum_jy_jP_j$, aplicar el lector a cada energía $Qy_j$
produce $\sum_jy_py_jP_j\otimes W_{\rm av}$, independientemente
de base propia. El cálculo espectral extiende el lector lineal
al carácter autoadjunto positivo general.
La descomposición por $P_A,P_V$ da los operadores
$y_pw_AY$ y $y_pw_VY$, con sus multiplicidades:
son autoadjuntos positivos en el dominio de $Y\otimes I$.
Sustituir los lectores radiales prueba las otras expresiones.
\end{proof}

Si una realización adicional identifica efectivamente su operador
completo con $H_X=QY$, conserva esa identificación y sus dominios.
Para $H_X=\sum_rH_r$ en un núcleo común,
$\widehat I_{p,Y}=(E_p/Q^2)\sum_rH_r\otimes W_{\rm av}$.
No se requiere positividad de cada sumando; sí la del carácter
del operador identificado. No se suprimen términos ni se cambia
el origen de energía para imponerla.
El teorema no asigna automáticamente este lector a $\mathbb H$.

En los bloques de \eqref{hmtcuatro:gea:schur-prueba}, el inverso
interior de $y_pY\otimes W_{\rm av}$ es
$y_p^{-1}C^{-1}\otimes W_{\rm av}^{-1}$.
La multiplicación da
\begin{equation}
\operatorname{Schur}(\widehat I_{p,Y})
=y_pY_{\rm ef}\otimes W_{\rm av}
=\frac{R_{X,\rm ef}}{\bar\lambda_p}\otimes W_{\rm av}.
\label{hmtcuatro:gea:schur-autoinercia}
\end{equation}
No se borran los bloques cruzados.
Para $a_s=y_pw_s>0$, $s=A,V$, la misma cuenta da
$\mathcal K_{a_sY}(a_sz)=a_s\mathcal K_Y(z)$
en el dominio resolvente interior, conservando la memoria completa.
Si $Y'K=KY$ y $W'_{\rm av}J=JW_{\rm av}$, con
$y_p$ y sección intensiva fijos,
\begin{equation}
\widehat I'_{p,Y'}(K\otimes J)
=(K\otimes J)\widehat I_{p,Y}.
\label{hmtcuatro:gea:naturalidad}
\end{equation}
Se prueba por sustitución.
Los cortes espectrales de la proposición
\ref{hmtcuatro:gea:limite} convergen en norma de grafo;
los dos pesos fijos transmiten la convergencia a la respuesta.
La naturalidad es la de estos entrelazadores efectivos, no la
de una familia distinta de mapas espaciales identificada sólo por nombre.

\subsection{Autoinercia del estado y aceleración de su proyección}
\label{hmtcuatro:gea:proyeccion}

En una realización diferenciable, el estado conjunto
$\Gamma=(x,\mathcal O,m)$ tiene conexión $\nabla^{\rm tot}$
en $\mathcal N$. Una proyección suave
$\pi_S:\mathcal N\to\mathcal N_S$, con conexión $\nabla^S$,
conserva los transportes y lectores que realiza.
El marco observado es
$\mathfrak F_{\mathcal O}(\Gamma)
=\operatorname{Res}_{\mathcal O}(\Gamma,\nabla^{\rm tot})$.
Para una curva dos veces diferenciable y
$\gamma_S=\pi_S\circ\Gamma$, la autoinercia significa
$\nabla^{\rm tot}_{\dot\Gamma}\dot\Gamma=0$.
El defecto es
\[
\mathcal K_{\pi_S}(\dot\Gamma,\dot\Gamma)
=\nabla^S_{\dot\gamma_S}\dot\gamma_S
-d\pi_S(\nabla^{\rm tot}_{\dot\Gamma}\dot\Gamma).
\]
\begin{proposition}[Aceleración publicada]
\label{hmtcuatro:gea:aceleracion}
Para una trayectoria autoinercial,
\begin{equation}
a_S=\mathcal K_{\pi_S}(\dot\Gamma,\dot\Gamma),\qquad
(\mathcal K_{\pi_S})^a{}_{BC}
=\partial_B\partial_C\pi_S^a
+(\Gamma^S)^a{}_{bc}\partial_B\pi_S^b\partial_C\pi_S^c
-\partial_D\pi_S^a(\Gamma^{\rm tot})^D{}_{BC}.
\label{hmtcuatro:gea:defecto}
\end{equation}
\end{proposition}
\begin{proof}
La regla de la cadena da
$\ddot\gamma_S^a=\partial_B\pi_S^a\ddot\Gamma^B
+\partial_B\partial_C\pi_S^a\dot\Gamma^B\dot\Gamma^C$.
Sumar la conexión del subsistema y restar la imagen de la
aceleración total cancela los términos en $\ddot\Gamma$.
Los restantes son el tensor mostrado contraído con las velocidades.
La autoinercia anula la aceleración total y deja la igualdad.
\end{proof}
La respuesta geométrica factoriza por $b$ cuando conexión,
proyección y dato cinemático descienden a la frontera, o éste
se fija en la comparación. Si cambia la velocidad, se conserva
como argumento. Así, el conjunto autoinercial puede publicar
subsistemas acelerados sin un marco absoluto exterior.

Autoinercia no anula por sí sola curvatura ni holonomía.
Un dato transportado por un lazo retorna si y sólo si está en
el espacio fijo de su holonomía. El retorno del estado completo
requiere además una realización inyectiva en el sector comparado:
retorno visible y retorno de historia siguen siendo distintos.

\subsection{Resultado y controles de alcance}
\label{hmtcuatro:gea:conclusion}

El mismo carácter fija masa, energía y radio; su proporcionalidad
conserva memoria estática, espectral y normalización dual.
Los transportes geométrico e interno entran en el diferencial
de una sola energía. La acción conjunta produce las fuentes
geométrica y torsional y conserva las restricciones sobre su
solución regular. El operador autoinercial es el radio del
carácter dividido por la longitud conjugada de la sonda,
mientras la aceleración es el defecto de proyección del estado.

Los dos últimos objetos tienen tipos distintos:
$\widehat I_{p,Y}$ no se identifica con $\mathcal K_{\pi_S}$
ni con el tensor de Einstein. La incorporación probada tampoco
se sustituye por la exigencia diferente de anular una curvatura
cuántica total para deformaciones normales arbitrarias.
Ese enunciado posee sus generadores y dominios; no se afirma aquí
ni se declara ausente del corpus por no ser esta conclusión.

Los falsadores son concretos. Cambiar sólo $G$ o $\hbar$ con
$L,Y$ fijos rompe \eqref{hmtcuatro:gea:radial}; sustituir
$L$ por $\ell_0$ sin $54/\pi$ altera el reloj; borrar $F$
modifica la reducción; reescalar juntos los lectores areal y
volumétrico no prueba separación machiana; un mapa que no
entrelace $Y$ no satisface \eqref{hmtcuatro:gea:naturalidad}.
Los ensayos racionales son controles finitos reproducibles,
no valores objetivo del generador ni sustitutos de las pruebas
de dominio y límite.
